% =============================================================================
% Zusammenfassung
% =============================================================================

\begin{abstract}
    \justifying
    Diese praxisorientierte Fallstudie analysiert die Verfügbarkeit und Fehlertoleranz
    des ROSI-Systems (\textit{Realtime Optical Surface Inspection}) im industriellen
    24/7-Betrieb. ROSI ist ein On-Premise-Inline-Scan-System der Grenzebach BSH GmbH
    zur automatisierten Qualitätskontrolle von Baustoffen und technisch als
    Python-Multiprocessing-System mit Shared-Memory-basierter Verarbeitungspipeline
    realisiert.

    Der Schwerpunkt liegt auf Fehlern, die im Dauerbetrieb schleichend entstehen, lange
    unbemerkt bleiben oder kaskadieren und im Entwicklungsbetrieb daher selten sichtbar
    werden. Fünf Failure Cases auf Prozess- und Betriebssystemebene wurden in einer
    kontrollierten Testumgebung mittels Fault Injection reproduzierbar untersucht.

    Die Messungen zeigen konkrete Lücken in der Fehlertoleranz: Drei von vier
    Service-Prozessen fallen ohne Erkennung aus; eine volle RAM-Disk führt nach
    35\,s Stillstand zum Absturz der Anlage; eine blockierende Log-Queue hält die
    Pipeline ohne Alarm an. Eine veraltete Datenbankverbindung bleibt in der
    Entwicklungsumgebung unsichtbar, ist in der Produktionskonfiguration aber als Risiko
    ableitbar. Eine angenommene Race Condition im Speichermanagement wurde widerlegt. Aus
    den Befunden werden aufwandsarme Maßnahmen zur besseren Überwachung und
    Fehlerentkopplung abgeleitet.
\end{abstract}
